\gdef\tagBrand{NT}
\gdef\footerAddress{Zalo: 0799666070}
\setcounter{ex}{0}

% --- CÂU 1 ---
\begin{ex}
	Cho bất phương trình bậc nhất hai ẩn $2x - 3y + 6 > 0$ và hai điểm $M(1; 2)$, $N(-4; 0)$. Trong các mệnh đề sau, mệnh đề nào đúng, mệnh đề nào sai?
	\choiceTFt
	{\True Bất phương trình đã cho là một bất phương trình bậc nhất hai ẩn $x$ và $y$}
	{\True Cặp số $(x; y) = (1; 2)$ là một nghiệm của bất phương trình đã cho}
	{Đường thẳng $d\colon 2x - 3y + 6 = 0$ đi qua gốc tọa độ $O(0; 0)$}
	{Điểm $N(-4; 0)$ thuộc miền nghiệm của bất phương trình đã cho}
	\loigiai{
		\begin{itemchoice}
			\itemch \textbf{Đúng.} Bất phương trình có dạng $ax + by + c > 0$ với $a = 2, b = -3$ ($a, b$ không đồng thời bằng $0$), do đó đây là bất phương trình bậc nhất hai ẩn.
			\itemch \textbf{Đúng.} Thay $(x; y) = (1; 2)$ vào vế trái bất phương trình, ta được:
			$$2 \cdot 1 - 3 \cdot 2 + 6 = 2 > 0\text{ (thỏa mãn)}.$$
			Vậy cặp số $(1; 2)$ là một nghiệm của bất phương trình.
			\itemch \textbf{Sai.} Thay tọa độ $O(0; 0)$ vào phương trình đường thẳng $d$, ta có:
			$$2 \cdot 0 - 3 \cdot 0 + 6 = 6 \ne 0.$$
			Do đó đường thẳng $d$ không đi qua gốc tọa độ $O(0; 0)$.
			\itemch \textbf{Sai.} Thay tọa độ $N(-4; 0)$ vào vế trái bất phương trình, ta được:
			$$2 \cdot (-4) - 3 \cdot 0 + 6 = -2 < 0.$$
			Do đó điểm $N(-4; 0)$ không thuộc miền nghiệm của bất phương trình.
		\end{itemchoice}
	}
\end{ex}

% --- CÂU 2 ---
\begin{ex}
	Cho góc $\alpha$ thỏa mãn $90^\circ < \alpha < 180^\circ$ và $\sin\alpha = \dfrac{3}{5}$. Trong các mệnh đề sau, mệnh đề nào đúng, mệnh đề nào sai?
	\choiceTFt
	{\True Giá trị của $\cos\alpha$ mang dấu âm}
	{\True $\cos\alpha = -\dfrac{4}{5}$}
	{$\tan\alpha = \dfrac{3}{4}$}
	{\True Giá trị của biểu thức $P = 2\sin(180^\circ - \alpha) + \cos(180^\circ - \alpha)$ bằng $2$}
	\loigiai{
		\begin{itemchoice}
			\itemch \textbf{Đúng.} Vì $90^\circ < \alpha < 180^\circ$ ($\alpha$ là góc tù) nên $\cos\alpha < 0$.
			\itemch \textbf{Đúng.} Ta có $\sin^2\alpha + \cos^2\alpha = 1 \implies \cos^2\alpha = 1 - \sin^2\alpha = 1 - \left(\dfrac{3}{5}\right)^2 = \dfrac{16}{25}$.\\
			Vì $\cos\alpha < 0$ nên $\cos\alpha = -\sqrt{\dfrac{16}{25}} = -\dfrac{4}{5}$.
			\itemch \textbf{Sai.} Ta có:
			$$\tan\alpha = \dfrac{\sin\alpha}{\cos\alpha} = \dfrac{\frac{3}{5}}{-\frac{4}{5}} = -\dfrac{3}{4}.$$
			\itemch \textbf{Đúng.} Với hai góc bù nhau:
			$$\sin(180^\circ - \alpha) = \sin\alpha = \dfrac{3}{5}\quad \text{và}\quad \cos(180^\circ - \alpha) = -\cos\alpha = -\left(-\dfrac{4}{5}\right) = \dfrac{4}{5}.$$
			Suy ra:
			$$P = 2 \cdot \dfrac{3}{5} + \dfrac{4}{5} = \dfrac{6}{5} + \dfrac{4}{5} = \dfrac{10}{5} = 2.$$
		\end{itemchoice}
	}
\end{ex}
